% ======================================================================
%  preamble.tex —— 全书共用导言区（main.tex 与各节自检文件共同载入）
%  只做版式与宏定义，不含正文；改版式只改这里。
%  底本：Lin Chen, NOTES FOR ALGEBRAIC GEOMETRY 1 (2025-12-27), 171 pp.
% ======================================================================
\usepackage[margin=1.0in]{geometry}
\parskip 1ex
\usepackage[no-math]{fontspec}
\usepackage{xeCJK}
\setCJKmainfont{SimSun}[BoldFont=SimHei,ItalicFont=KaiTi,BoldItalicFont=SimHei]
\setCJKsansfont{SimHei}[BoldFont=SimHei]
\setCJKmonofont{SimSun}
\setCJKfamilyfont{kaishu}{KaiTi}
\newcommand{\kaishu}{\CJKfamily{kaishu}}
\setmainfont{lmroman10-regular.otf}[BoldFont=lmroman10-bold.otf,ItalicFont=lmroman10-italic.otf,BoldItalicFont=lmroman10-bolditalic.otf]
\setsansfont{lmsans10-regular.otf}[BoldFont=lmsans10-bold.otf,ItalicFont=lmsans10-oblique.otf,BoldItalicFont=lmsans10-boldoblique.otf]
\setmonofont{lmmono10-regular.otf}

\usepackage{amsmath,amssymb,amsthm,mathtools,bm}
\usepackage{mathrsfs}
\usepackage{aliascnt}
\usepackage{graphicx}
\usepackage{array}
\usepackage{longtable}
\usepackage{enumitem}
\usepackage{booktabs}
\usepackage{float}
\usepackage[font=normalsize,labelsep=period]{caption}
\usepackage[export]{adjustbox}
\usepackage{tikz-cd}
\usepackage{titlesec}
\titleformat*{\section}{\normalfont\sffamily\Large\bfseries}
\titleformat*{\subsection}{\normalfont\sffamily\large\bfseries}
\titleformat*{\subsubsection}{\normalfont\sffamily\normalsize\bfseries}
\usepackage[section]{placeins}
\usepackage{xcolor}
\usepackage{framed}
\usepackage[unicode,hidelinks,bookmarksnumbered]{hyperref}
\usepackage{bookmark}
\usepackage[nameinlink,noabbrev]{cleveref}

\hypersetup{
  pdftitle={代数几何讲义 I（中译注本）},
  pdfauthor={Lin Chen 著；中译注},
  pdfsubject={代数几何，概形，层论与上同调（底本 2025-12-27 版，171 页）},
  pdfkeywords={代数几何, 概形, 层, 拟凝聚层, 射影态射, 中译注本}
}

\renewcommand{\contentsname}{目录}
\renewcommand{\refname}{参考文献}
\renewcommand{\figurename}{图}
\renewcommand{\tablename}{表}
\renewcommand{\proofname}{\bfseries 证明}
\renewcommand{\partname}{第}
\renewcommand{\abstractname}{摘要}
\renewcommand{\appendixname}{附录}
\renewcommand{\indexname}{索引}

\pagestyle{plain}
\setcounter{secnumdepth}{3}
\setcounter{tocdepth}{2}

% ---------- 版式：行距与页眉页脚 ----------
\linespread{1.12}
\usepackage{fancyhdr}
\setlength{\headheight}{14pt}
\pagestyle{fancy}
\fancyhf{}
\fancyhead[L]{\small\kaishu 代数几何讲义 I\,·\,中译注本}
\fancyhead[R]{\small\kaishu\nouppercase{\rightmark}}
\fancyfoot[C]{\small\thepage}
\renewcommand{\headrulewidth}{0.4pt}
\renewcommand{\footrulewidth}{0pt}
\fancypagestyle{plain}{\fancyhf{}\fancyfoot[C]{\small\thepage}\renewcommand{\headrulewidth}{0pt}}
\renewcommand{\sectionmark}[1]{\markright{\thesection\quad #1}}
\renewcommand{\subsectionmark}[1]{}

% ---------- 译注环境（全局开关 \ifshowyizhu） ----------
% 用法：\begin{yizhu}……[来源：……]\end{yizhu}
\newif\ifshowyizhu
\showyizhutrue
\newenvironment{yizhubox}{%
  \def\FrameCommand{{\color{black!55}\vrule width 1.4pt}\hspace{9pt}}%
  \MakeFramed{\advance\hsize-\width\sloppy\FrameRestore}}%
 {\endMakeFramed}
\newenvironment{yizhu}{%
  \ifshowyizhu\begin{yizhubox}\small\noindent{\bfseries 译注}\quad\ignorespaces\fi}%
 {\ifshowyizhu\end{yizhubox}\fi}

% ---------- 定理类环境：统一计数器，编号 = 节.小节.序号（同原书） ----------
\theoremstyle{plain}
\newtheorem{theorem}{定理}[subsection]
\newaliascnt{lemma}{theorem}\newtheorem{lemma}[lemma]{引理}\aliascntresetthe{lemma}
\newaliascnt{proposition}{theorem}\newtheorem{proposition}[proposition]{命题}\aliascntresetthe{proposition}
\newaliascnt{corollary}{theorem}\newtheorem{corollary}[corollary]{推论}\aliascntresetthe{corollary}
\newaliascnt{claim}{theorem}\newtheorem{claim}[claim]{断言}\aliascntresetthe{claim}
\newaliascnt{conjecture}{theorem}\newtheorem{conjecture}[conjecture]{猜想}\aliascntresetthe{conjecture}

\theoremstyle{definition}
\newaliascnt{definition}{theorem}\newtheorem{definition}[definition]{定义}\aliascntresetthe{definition}
\newaliascnt{example}{theorem}\newtheorem{example}[example]{例}\aliascntresetthe{example}
\newaliascnt{predefinition}{theorem}\newtheorem{predefinition}[predefinition]{预备定义}\aliascntresetthe{predefinition}
\newaliascnt{remark}{theorem}\newtheorem{remark}[remark]{注记}\aliascntresetthe{remark}
\newaliascnt{notation}{theorem}\newtheorem{notation}[notation]{记号}\aliascntresetthe{notation}
\newaliascnt{construction}{theorem}\newtheorem{construction}[construction]{构造}\aliascntresetthe{construction}
\newaliascnt{exercise}{theorem}\newtheorem{exercise}[exercise]{习题}\aliascntresetthe{exercise}
\newaliascnt{warning}{theorem}\newtheorem{warning}[warning]{警告}\aliascntresetthe{warning}
\newaliascnt{question}{theorem}\newtheorem{question}[question]{问题}\aliascntresetthe{question}
\newaliascnt{convention}{theorem}\newtheorem{convention}[convention]{约定}\aliascntresetthe{convention}

\newenvironment{sketch}[1][证明概略]{\begin{proof}[#1]}{\end{proof}}
\newenvironment{proofsketch}[1][证明概略]{\begin{proof}[#1]}{\end{proof}}
\newenvironment{proofidea}[1][证明思路]{\begin{proof}[#1]}{\end{proof}}

\crefname{theorem}{定理}{定理}
\crefname{lemma}{引理}{引理}
\crefname{proposition}{命题}{命题}
\crefname{corollary}{推论}{推论}
\crefname{claim}{断言}{断言}
\crefname{conjecture}{猜想}{猜想}
\crefname{definition}{定义}{定义}
\crefname{example}{例}{例}
\crefname{predefinition}{预备定义}{预备定义}
\crefname{remark}{注记}{注记}
\crefname{notation}{记号}{记号}
\crefname{construction}{构造}{构造}
\crefname{exercise}{习题}{习题}
\crefname{warning}{警告}{警告}
\crefname{question}{问题}{问题}
\crefname{convention}{约定}{约定}
\crefname{equation}{式}{式}
\crefname{figure}{图}{图}
\crefname{section}{节}{节}
\crefname{subsection}{小节}{小节}

% ---------- 数学宏 ----------
\newcommand{\Z}{\mathbb{Z}}
\newcommand{\Q}{\mathbb{Q}}
\newcommand{\R}{\mathbb{R}}
\newcommand{\C}{\mathbb{C}}
\newcommand{\N}{\mathbb{N}}
\newcommand{\F}{\mathbb{F}}
\newcommand{\A}{\mathbb{A}}
\newcommand{\PP}{\mathbb{P}}
\newcommand{\GG}{\mathbb{G}}
\newcommand{\bbk}{\mathbb{k}}
\newcommand{\cO}{\mathcal{O}}
\newcommand{\cM}{\mathcal{M}}
\newcommand{\cN}{\mathcal{N}}
\newcommand{\cF}{\mathcal{F}}
\newcommand{\cG}{\mathcal{G}}
\newcommand{\cI}{\mathcal{I}}
\newcommand{\cJ}{\mathcal{J}}
\newcommand{\cL}{\mathcal{L}}
\newcommand{\cE}{\mathcal{E}}
\newcommand{\cA}{\mathcal{A}}
\newcommand{\cB}{\mathcal{B}}
\newcommand{\cC}{\mathcal{C}}
\newcommand{\cP}{\mathcal{P}}
\newcommand{\cS}{\mathcal{S}}
\newcommand{\cU}{\mathcal{U}}
\newcommand{\cD}{\mathcal{D}}
\newcommand{\cX}{\mathcal{X}}
\newcommand{\cY}{\mathcal{Y}}
\newcommand{\Spec}{\operatorname{Spec}}
\newcommand{\Spm}{\operatorname{Spm}}
\newcommand{\Smax}{\operatorname{Spm}}
\newcommand{\Proj}{\operatorname{Proj}}
\newcommand{\Spv}{\operatorname{Spv}}
\newcommand{\Hom}{\operatorname{Hom}}
\newcommand{\uHom}{\underline{\mathrm{Hom}}}
\newcommand{\End}{\operatorname{End}}
\newcommand{\Aut}{\operatorname{Aut}}
\newcommand{\Isom}{\operatorname{Isom}}
\newcommand{\id}{\operatorname{id}}
\newcommand{\coker}{\operatorname{coker}}
\newcommand{\im}{\operatorname{im}}
\newcommand{\Ker}{\operatorname{Ker}}
\newcommand{\Ob}{\operatorname{Ob}}
\newcommand{\Mor}{\operatorname{Mor}}
\newcommand{\Frac}{\operatorname{Frac}}
\newcommand{\Pic}{\operatorname{Pic}}
\newcommand{\Supp}{\operatorname{Supp}}
\newcommand{\cHom}{\underline{\mathcal{H}om}}
\newcommand{\Ext}{\operatorname{Ext}}
\newcommand{\Tor}{\operatorname{Tor}}
\newcommand{\colim}{\operatorname*{colim}}
\newcommand{\CRing}{\mathsf{CRing}}
\newcommand{\Ring}{\mathsf{Ring}}
\newcommand{\Set}{\mathsf{Set}}
\newcommand{\Ab}{\mathsf{Ab}}
\newcommand{\Grp}{\mathsf{Grp}}
\newcommand{\Top}{\mathsf{Top}}
\newcommand{\Sch}{\mathsf{Sch}}
\newcommand{\PShv}{\mathsf{PShv}}
\newcommand{\Shv}{\mathsf{Shv}}
\newcommand{\Mod}{\mathsf{Mod}}
\newcommand{\Vect}{\mathsf{Vect}}
\newcommand{\op}{\mathrm{op}}
\newcommand{\smallcat}[1]{\mathsf{#1}}
\newcommand{\omod}{\cO_X\text{-}\mathsf{mod}}
\newcommand{\omodP}{\cO_X\text{-}\mathsf{mod}^{\PShv}}
\newcommand{\lto}{\longrightarrow}
\newcommand{\lfrom}{\longleftarrow}
\newcommand{\into}{\hookrightarrow}
\newcommand{\from}{\leftarrow}
\newcommand{\eop}{\unskip\nobreak\hfill\ensuremath{\square}}
\providecommand{\longmapsfrom}{\mathrel{\rotatebox[origin=c]{180}{$\longmapsto$}}}
\providecommand{\mapsfrom}{\mathrel{\rotatebox[origin=c]{180}{$\mapsto$}}}
\newcommand{\secref}[1]{第~\ref{sec:#1}~节}
\newcommand{\secrefx}[1]{\ref{sec:#1}}
\pdfstringdefDisableCommands{\def\PShv{PShv}\def\Shv{Shv}\def\Spec{Spec}\def\Proj{Proj}\def\cO{O}\def\Z{Z}\def\A{A}\def\PP{P}\def\omod{mod}}

% ---- 通用补充宏（各节执笔者可能用到，统一在此定义，避免 undefined） ----
\providecommand{\supp}{\operatorname{Supp}}
\providecommand{\Ann}{\operatorname{Ann}}
\providecommand{\ann}{\operatorname{ann}}
\providecommand{\Id}{\mathrm{Id}}
\providecommand{\rank}{\operatorname{rank}}
\providecommand{\tr}{\operatorname{tr}}
\providecommand{\coker}{\operatorname{coker}}
\providecommand{\cH}{\mathcal{H}}
\providecommand{\cK}{\mathcal{K}}
\providecommand{\cQ}{\mathcal{Q}}
\providecommand{\cR}{\mathcal{R}}
\providecommand{\cT}{\mathcal{T}}
\providecommand{\cV}{\mathcal{V}}
\providecommand{\cW}{\mathcal{W}}
\providecommand{\cZ}{\mathcal{Z}}
\providecommand{\cK}{\mathcal{K}}
\providecommand{\cJ}{\mathcal{J}}
\providecommand{\cG}{\mathcal{G}}
\providecommand{\cS}{\mathcal{S}}
\providecommand{\cU}{\mathcal{U}}
\providecommand{\cX}{\mathcal{X}}
\providecommand{\cY}{\mathcal{Y}}

\providecommand{\coim}{\operatorname{coim}}
\providecommand{\QCoh}{\mathsf{QCoh}}
\providecommand{\Coh}{\mathsf{Coh}}
\providecommand{\Sym}{\operatorname{Sym}}
\makeatletter
\@ifundefined{propdef}{\newtheorem{propdef}[theorem]{命题—定义}\crefname{propdef}{命题—定义}{命题—定义}}{}
\makeatother

\counterwithin*{equation}{section}
\renewcommand{\theequation}{\thesection.\arabic{equation}}
